% (c) Nikita Lisitsa, lisyarus@gmail.com, 2026

\documentclass[10pt]{beamer}

\usepackage[T2A]{fontenc}
\usepackage[russian]{babel}
\usepackage{minted}

\usepackage[most]{tcolorbox}
\tcbuselibrary{minted}

\usepackage{graphicx}
\graphicspath{ {./images/} }

\usepackage{adjustbox}

\usepackage{color}
\usepackage{soul}
\usepackage{multicol}
\usepackage{multirow}
\usepackage{hyperref}
\usepackage{tikz}
\usetikzlibrary{arrows.meta,positioning,calc}

\usetheme{metropolis}

\definecolor{red}{rgb}{1,0,0}
\definecolor{codebg}{RGB}{29,35,49}
\setminted{fontsize=\footnotesize}

\makeatletter
\newcommand{\slideimage}[1]{
  \begin{figure}
    \begin{adjustbox}{width=\textwidth, totalheight=\textheight-2\baselineskip-2\baselineskip,keepaspectratio}
      \includegraphics{#1}
    \end{adjustbox}
  \end{figure}
}
\makeatother

\title{Язык программирования C}
\subtitle{Лекция 4: динамическая память (куча), структуры}
\author{Лисица Никита}
\date{2026}

\setbeamertemplate{footline}[frame number]

\usemintedstyle{tango}

\begin{document}

\frame{\titlepage}

\begin{frame}[fragile]
\frametitle{Напоминание: три вида памяти}
\begin{itemize}
\item \textbf{Стек}: выделяется на старте программы, размер предопределён (стек фреймы ``выделяет'' на стеке компилятор)
\pause
\item \textbf{Глобальная память}: выделяется на старте программы, размер известен на момент компиляции (память, нужная под все глобальные переменные)
\pause
\item \textbf{Динамическая память (куча)}: выделяется программистом, размер произвольный
\end{itemize}
\end{frame}

\begin{frame}[fragile]
\frametitle{Динамическая память (куча)}
\begin{itemize}
\item \mintinline{cpp}|void* malloc(size_t size);| -- выделяет \underline{непрерывный} кусок памяти размером \mintinline{cpp}|size| байт
\pause
\item \mintinline{cpp}|void free(void *ptr);| -- возвращает память системе для переиспользования
\pause
\item Если не удалось выделить память (напр. закончилась свободная память в системе), \mintinline{cpp}|malloc| возвращает \textit{нулевой указатель}: \mintinline{cpp}|NULL|
\end{itemize}
\end{frame}

\begin{frame}[fragile]
\frametitle{Динамическая память (куча)}
\usemintedstyle{lightbulb}
\begin{tcblisting}{listing engine=minted, minted language=cpp, minted options={fontsize=\scriptsize}, listing only, colback=codebg, colframe=codebg, rounded corners}
// Выделить память под миллион объектов типа int
int *p = malloc(1000000 * sizeof(int));

// То же самое: if (!p)
if (p == NULL) {
  // Обработаем нехватку памяти
}

// Можем использовать p как указатель на массив:
for (size_t i = 0; i < 1000000; ++i) {
  p[i] = i * i;
}

// Освобождаем память
free(p);
\end{tcblisting}
\end{frame}

\begin{frame}[fragile]
\frametitle{Динамическая память (куча)}
\begin{itemize}
\item Размер выделяемой памяти может быть полностью динамическим (введён пользователем, вычислен на основе других данных, и т.п.) -- невозможно для стека и глобальных переменных
\end{itemize}
\usemintedstyle{lightbulb}
\begin{tcblisting}{listing engine=minted, minted language=cpp, minted options={fontsize=\scriptsize}, listing only, colback=codebg, colframe=codebg, rounded corners}
size_t size;
scanf("%zu", &size);

char *p = malloc(size);

// Как-то используем p...

free(p);
\end{tcblisting}
\end{frame}

\begin{frame}[fragile]
\frametitle{Динамическая память (куча)}
\begin{itemize}
\item Как правило, маленькие объекты нет смысла хранить в куче -- их дешевле скопировать
\pause
\item \mintinline{cpp}|int *p = malloc(sizeof(int));| -- тратим 4 байта на сам int (в куче) и 8 байт на указатель (на стеке)
\pause
\begin{itemize}
\item И ещё 8-16 байт в куче для метаданных аллокатора :)
\end{itemize}
\pause
\item C и C++ позволяют явно управлять расположением объекта в памяти -- этим стоит пользоваться :)
\end{itemize}
\end{frame}

\begin{frame}[fragile]
\frametitle{Утечки памяти}
\begin{multicols}{2}
\usemintedstyle{lightbulb}
\begin{tcblisting}{listing engine=minted, minted language=cpp, minted options={fontsize=\scriptsize}, listing only, colback=codebg, colframe=codebg, rounded corners}
int foo(size_t size) {
  int *p = malloc(size);
  // Что-то посчитали
  // Забыли вызвать free
  return p[0];
}
\end{tcblisting}
\columnbreak
\begin{tcblisting}{listing engine=minted, minted language=cpp, minted options={fontsize=\scriptsize}, listing only, colback=codebg, colframe=codebg, rounded corners}
int *p = malloc(size);
// Что-то посчитали
// Забыли вызвать free

// Перевыделили память
p = malloc(new_size);
\end{tcblisting}
\end{multicols}
\usemintedstyle{tango}
\begin{itemize}
\item \textit{Утечка памяти}: мы ``потеряли'' значение указателя на выделенную область \(\Longrightarrow\) не можем её освободить (нам не от чего вызвать \mintinline{cpp}|free|)
\end{itemize}
\end{frame}

\begin{frame}[fragile]
\frametitle{Утечки памяти}
\begin{itemize}
\item При завершении программы память автоматически вернётся системе, зачем её освобождать?
\pause
\item Многие программы работают продолжительное время (браузер, сервер, игра), и память действительно может закончиться
\pause
\item Код, который вы пишете для ``одноразовой'' программы, потом могут переиспользовать для другого проекта
\pause
\item Для поиска утечек есть много инструментов: \textbf{Valgrind}, \textbf{Address Sanitizer}
\end{itemize}
\end{frame}

\begin{frame}[fragile]
\frametitle{Динамическая память (куча)}
\begin{itemize}
\item \mintinline{cpp}|malloc|/\mintinline{cpp}|free| реализованы в стандартной библиотеке C
\pause
\item Они используют специальные потокобезопасные структуры данных, чтобы
\begin{itemize}
\item Выделять новые участки памяти
\item Переиспользовать старые
\item Выполнять дефрагментацию (склеивать соседние неиспользуемые участки)
\end{itemize}
\pause
\item \(\Longrightarrow\) Выделение динамической памяти занимает некоторое время (сотни наносекунд / десятки микросекунд)
\pause
\item \(\Longrightarrow\) Лучше выделять и переиспользовать большой кусок памяти, чем выделять много маленьких объектов
\end{itemize}
\end{frame}

\begin{frame}[fragile]
\frametitle{Двумерный массив на куче: вариант 1}
\usemintedstyle{lightbulb}
\begin{tcblisting}{listing engine=minted, minted language=cpp, minted options={fontsize=\scriptsize}, listing only, colback=codebg, colframe=codebg, rounded corners}
int **m = malloc(N * sizeof(int*));
for (int i = 0; i < N; ++i)
  m[i] = malloc(N * sizeof(int));

// Можем работать с m[i][j]

// Очищать придётся циклом
for (int i = 0; i < N; ++i)
  free(m[i]);
free(m);
\end{tcblisting}
\end{frame}

\begin{frame}[fragile]
\frametitle{Двумерный массив на куче: вариант 2}
\usemintedstyle{lightbulb}
\begin{tcblisting}{listing engine=minted, minted language=cpp, minted options={fontsize=\scriptsize}, listing only, colback=codebg, colframe=codebg, rounded corners}
int **m = malloc(N * sizeof(int*));
m[0] = malloc(N * N * sizeof(int));
for (int i = 1; i < N; ++i)
  m[i] = m[0] + i * N;

// Можем работать с m[i][j]

free(m[0]);
free(m);
\end{tcblisting}
\end{frame}

\begin{frame}[fragile]
\frametitle{Двумерный массив на куче: вариант 3}
\usemintedstyle{lightbulb}
\begin{tcblisting}{listing engine=minted, minted language=cpp, minted options={fontsize=\scriptsize}, listing only, colback=codebg, colframe=codebg, rounded corners}
int *m = malloc(N * N * sizeof(int));

// Для обращения к элементу (i,j) придётся явно вычислить его адрес
// Обычно это быстрее, чем прочитать адрес строки m[i] из памяти
m[i * N + j] = 42;

free(m);
\end{tcblisting}
\end{frame}

\begin{frame}[fragile]
\frametitle{Динамическая память (куча)}
\begin{itemize}
\item \mintinline{cpp}|calloc| -- выделить память и заполнить нулями (нулевыми байтами)
\pause
\item \mintinline{cpp}|realloc| -- \textit{попробовать} увеличить размер уже выделенного куска
\pause
\begin{itemize}
\item Если после него есть свободная память, увеличит размер и вернёт тот же указатель
\pause
\item Если после него нет свободной памяти, выделит новый кусок нужного размера, скопирует туда старые данные, и освободит старый кусок
\end{itemize}
\end{itemize}
\end{frame}

\begin{frame}[fragile]
\frametitle{Структуры}
\begin{itemize}
\item Часто хочется объединить несколько переменных в одну логическую единицу:
\pause
\begin{itemize}
\item Точка в пространстве: 3 числа \mintinline{cpp}|float x, y, z;|
\pause
\item Ребро взвешенного графа: \mintinline{cpp}|int start, end; float weight;|
\pause
\item Вершина бинарного дерева: \mintinline{cpp}|node *left, *right; int data;|
\pause
\item Состояние тайла 2D игры: \mintinline{cpp}|int terrain; bool has_road;|
\end{itemize}
\pause
\item В языке C для этого есть встроенный механизм: \textit{структуры}
\end{itemize}
\end{frame}

\begin{frame}[fragile]
\frametitle{Структуры}
\usemintedstyle{lightbulb}
\begin{multicols}{2}
\begin{tcblisting}{listing engine=minted, minted language=cpp, minted options={fontsize=\scriptsize}, listing only, colback=codebg, colframe=codebg, rounded corners}
struct point_3d {
  float x, y, z;
};
\end{tcblisting}
\begin{tcblisting}{listing engine=minted, minted language=cpp, minted options={fontsize=\scriptsize}, listing only, colback=codebg, colframe=codebg, rounded corners}
struct node {
  struct node *left;
  struct node *right;
  int data;
}
\end{tcblisting}
\columnbreak
\begin{tcblisting}{listing engine=minted, minted language=cpp, minted options={fontsize=\scriptsize}, listing only, colback=codebg, colframe=codebg, rounded corners}
struct edge {
  int start, end;
  float weight;
};
\end{tcblisting}
\begin{tcblisting}{listing engine=minted, minted language=cpp, minted options={fontsize=\scriptsize}, listing only, colback=codebg, colframe=codebg, rounded corners}
struct tile {
  int terrain;
  bool has_road;
};
\end{tcblisting}
\end{multicols}
\end{frame}

\begin{frame}[fragile]
\frametitle{Структуры}
\usemintedstyle{lightbulb}
\begin{tcblisting}{listing engine=minted, minted language=cpp, minted options={fontsize=\scriptsize}, listing only, colback=codebg, colframe=codebg, rounded corners}
// Объявить переменную типа структуры
struct edge e;

// Прочитать поля структуры из файла
scanf("%d %d %f", &e.start, &e.end, &e.weight);

// Завести массив структур
struct edge e[10];
e[4].start = 15;

// Выделить динамический массив структур
struct edge *e = malloc(100 * sizeof(struct edge));

// Обращение к полю через указатель
// (эквивалентные способы)
e[0].weight = 1.412f;
(*e).weight = 3.1415f;
e->weight = 2.71828f;
\end{tcblisting}
\end{frame}

\begin{frame}[fragile]
\frametitle{Структуры: инициализация}
\usemintedstyle{lightbulb}
\begin{tcblisting}{listing engine=minted, minted language=cpp, minted options={fontsize=\scriptsize}, listing only, colback=codebg, colframe=codebg, rounded corners}
// Инициализация:
struct edge e = { 10, 15, 1.2345f };

// Designated инициализация:
struct edge e = {
  .start = 10,
  .end = 15,
  .weight = 1.2345f,
};
\end{tcblisting}
\end{frame}

\begin{frame}[fragile]
\frametitle{Структуры: копирование}
\usemintedstyle{lightbulb}
\begin{tcblisting}{listing engine=minted, minted language=cpp, minted options={fontsize=\scriptsize}, listing only, colback=codebg, colframe=codebg, rounded corners}
struct edge e1 = { 10, 15, 1.2345f };

// Полностью копирует все поля
struct edge e2 = e1;

// Тоже копирует все поля
e1 = e2;

e1.weight *= 2.f;

// Выведет разные значения
printf("%f %f", e1.weight, e2.weight);
\end{tcblisting}
\end{frame}

\begin{frame}[fragile]
\frametitle{Структуры: копирование}
\usemintedstyle{lightbulb}
\begin{itemize}
\item Если в структуре хранится указатель, то скопируется само \textit{значение указателя}, а не данные, на которые он указывает
\end{itemize}
\begin{tcblisting}{listing engine=minted, minted language=cpp, minted options={fontsize=\scriptsize}, listing only, colback=codebg, colframe=codebg, rounded corners}
struct array {
  int *data;
  int size;
};

struct array a = {
  .data = malloc(100 * sizeof(int)),
  .size = 100,
};

// b.data указывает на ту же память!
struct array b = a;

b.data[0] = 42;
printf("%i", a.data[0]); // выведет 42
\end{tcblisting}
\end{frame}

\begin{frame}[fragile]
\frametitle{Структуры: память}
\begin{itemize}
\item Как и все объекты, структуры в C можно хранить на стеке (локальная переменная), в глобальной памяти (глобальная/статическая переменная), или на куче (выделена через \mintinline{cpp}|malloc|)
\pause
\item Во многих других языках (Java, C\#, Python) составные объекты могут храниться только в куче
\end{itemize}
\end{frame}

\begin{frame}[fragile]
\frametitle{Структуры и функции}
\usemintedstyle{lightbulb}
\begin{tcblisting}{listing engine=minted, minted language=cpp, minted options={fontsize=\scriptsize}, listing only, colback=codebg, colframe=codebg, rounded corners}
struct big_struct {
  char name[1024];
  int values[256];
};

struct big_struct foo(struct big_struct value) {
  // ...
}
\end{tcblisting}
\pause
\begin{itemize}
\item Для передачи структуры в фунцию придётся выделить на стеке место и скопировать туда аргумент функции
\pause
\item Для возвращения структуры из фунции придётся сделать то же самое
\end{itemize}
\end{frame}

\begin{frame}[fragile]
\frametitle{Структуры и функции}
\begin{itemize}
\item В таких случаях лучше передавать указатель на структуру:
\end{itemize}
\usemintedstyle{lightbulb}
\begin{tcblisting}{listing engine=minted, minted language=cpp, minted options={fontsize=\scriptsize}, listing only, colback=codebg, colframe=codebg, rounded corners}
struct big_struct {
  char name[1024];
  int values[256];
};

void foo(struct big_struct *ptr) {
  // ...
}
\end{tcblisting}
\end{frame}

\begin{frame}[fragile]
\frametitle{Структуры: typedef}
\begin{itemize}
\item Писать каждый раз \mintinline{cpp}|struct my_struct| надоедает :)
\pause
\item Вместо этого можно сделать \mintinline{cpp}|typedef| -- синоним типа для структуры
\end{itemize}
\usemintedstyle{lightbulb}
\begin{tcblisting}{listing engine=minted, minted language=cpp, minted options={fontsize=\scriptsize}, listing only, colback=codebg, colframe=codebg, rounded corners}
// Новое имя для типа int
typedef int vertex_id;

// Новое имя для типа указателя на функцию
typedef int (*my_func_type)(int, int);

struct my_struct_t {
  char name[64];
  float weight;
};

// Новое имя для типа структуры
typedef struct my_struct my_struct;

// Теперь можно не писать struct!
my_struct x = { "Hello", 6.7f };
\end{tcblisting}
\end{frame}

\begin{frame}[fragile]
\frametitle{Структуры: typedef}
\begin{itemize}
\item \mintinline{cpp}|typedef| и объявление структуры можно объединить в одно объявление:
\end{itemize}
\usemintedstyle{lightbulb}
\begin{tcblisting}{listing engine=minted, minted language=cpp, minted options={fontsize=\scriptsize}, listing only, colback=codebg, colframe=codebg, rounded corners}
typedef struct my_struct {
  char name[64];
  float weight;
} my_struct;

// Или даже так:
typedef struct {
  char name[64];
  float weight;
} my_struct;
\end{tcblisting}
\end{frame}

\begin{frame}[fragile]
\frametitle{Структуры: typedef}
\begin{itemize}
\item Не стоит называть структуры в духе \mintinline{cpp}|my_struct_t|: стандарт POSIX резервирует все имена, заканчивающиеся на \mintinline{cpp}|_t| -- можем случайно пересечься с именами под Linux/macOS
\pause
\item Чем \mintinline{cpp}|typedef| лучше, чем директива препроцессора \mintinline{cpp}|#define|? \pause \mintinline{cpp}|#define| применяется ко всем идентификаторам, а \mintinline{cpp}|typedef| -- только к \textit{типам}
\end{itemize}
\end{frame}

\begin{frame}[fragile]
\frametitle{Структуры: выравнивание}
\usemintedstyle{lightbulb}
\begin{tcblisting}{listing engine=minted, minted language=cpp, minted options={fontsize=\scriptsize}, listing only, colback=codebg, colframe=codebg, rounded corners}
struct test {
  int i;
  char c;
};
\end{tcblisting}
\pause
\usemintedstyle{tango}
\begin{itemize}
\item Чему будет равно \mintinline{cpp}|sizeof(struct test)|? \pause На большинстве платформ -- 8:
\begin{itemize}
\item 4 байта на \mintinline{cpp}|int i|
\item 1 байт на \mintinline{cpp}|char c|
\item 3 байта на \textit{padding}
\end{itemize}
\end{itemize}
\end{frame}

\begin{frame}[fragile]
\frametitle{Структуры: выравнивание}
\begin{itemize}
\item Компилятор может сгенерировать более эффективный код, если адрес объекта типа \mintinline{cpp}|int| кратен 4 байтам (``выровнен'' на 4 байта)
\pause
\begin{itemize}
\item Обычно, выравнивание для примитивного типа равно его размеру
\end{itemize}
\pause
\item \(\Longrightarrow\) По умолчанию компилятор вставляет дополнительные неиспользуемые байты (\textit{padding}), чтобы обеспечить это
\end{itemize}
\end{frame}

\end{document}
